\section{Latihan}
\label{sec:latihan-bab3}

\begin{enumerate}
    \item \textbf{Konseptual.} Jelaskan perbedaan bahasa mesin, bahasa
          assembly, dan bahasa tingkat tinggi. Sertakan satu contoh
          operasi sederhana yang ditulis pada ketiga tingkat tersebut.

    \item \textbf{Konseptual.} Mengapa hubungan antara bahasa assembly
          dan bahasa mesin disebut \eng{one-to-one}? Bandingkan dengan
          hubungan bahasa tingkat tinggi dan bahasa mesin.

    \item \textbf{Konversi.} Ubah 156 desimal menjadi biner, oktal,
          dan heksadesimal. Tuliskan langkah pembagian atau
          pengelompokan bit yang Anda pakai.

    \item \textbf{Konversi dan tafsir.} Diberikan \bin{10110110}.
          Tentukan nilai desimalnya sebagai bilangan tidak bertanda
          dan sebagai bilangan bertanda 8 bit. Jelaskan perbedaan
          hasilnya.

    \item \textbf{Two's complement.} Tentukan bentuk 8 bit dari $-45$
          dengan langkah membalik bit dan menambahkan 1. Periksa
          hasilnya dengan mengurangkan 256 dari nilai tidak
          bertandanya.

    \item \textbf{Literal assembler.} Jelaskan mengapa \code{0F3H}
          ditulis dengan nol di depan, sedangkan \code{F3H} dapat
          menimbulkan kesalahan. Sebutkan pula dua kesalahan lain
          dalam penulisan literal bilangan.

    \item \textbf{Penulisan data.} Tulislah potongan segmen data yang
          mendefinisikan nilai 156 dalam bentuk desimal, biner, dan
          heksadesimal, lalu satu byte karakter \code{'Z'}.
          Prediksikan empat byte yang akan tampil pada DEBUG.

    \item \textbf{Analisis.} Sebuah pemeriksaan memori menghasilkan
          \code{C8 C8 C8 5A}. Jelaskan makna setiap byte, baik sebagai
          bilangan tidak bertanda maupun sebagai karakter.
\end{enumerate}

\begin{refleksibox}
\vspace{1.1em}
Jawab pertanyaan berikut pada jurnal praktikum Anda.
\begin{enumerate}
    \item Bagian konversi mana yang paling sulit Anda kerjakan?
          Langkah apa yang membantu Anda menyelesaikannya?
    \item Setelah memeriksa memori dengan DEBUG, apakah pemahaman Anda
          tentang hubungan bilangan dan data berubah? Jelaskan.
    \item Konsep apa yang masih perlu Anda ulangi sebelum mempelajari
          arsitektur 8086 pada bab berikutnya?
\end{enumerate}
\end{refleksibox}
